\documentclass[11pt, a4paper]{article}
\usepackage[margin=1in, top=1.2in, bottom=1.2in]{geometry}
% ---- Font & Micro-Typography ----
\usepackage[T1]{fontenc}
\usepackage{amsmath}
\usepackage{lmodern}
\usepackage{microtype}
% ---- Formatting & Aesthetics ----
\usepackage{graphicx}
% ---- Bibliography ----
\usepackage[backend=biber, style=numeric-comp, sorting=nyt, giveninits=true]{biblatex}


% ====================================================================================
\begin{document}

\begin{center}
  \LARGE\textbf{Daily Rant Digital Product Creation}\\
  Tadouanla Guetchuin Billy\\
  28177
\end{center}

Today in the center of knowledge \textbf{Digital Product Creation}, I had the opportunity to have
learn about 3D scanning. \textbf{3D scanning} is the process of capturing real-world objects
environments to be remodelled or analysed in the digital world. 3D scanning could also be seen as
reverse engineering process. For example: in the case where access to a particular part, spare
part, or even damage components, 3D scanning could easily be of used in such cases without
actually having access to the official dimensions from the manufacturer.

During the lab session we looked into the various scanners and utilizing different scanning
technologies and methods. We used a hand scanner, which uses a set of two cameras and laser
stripes. In order to scan an object, the object needs to be marked with some special markers as
well as the surface. These are used as reference point for the scanner in order for it to locate
itself. The marked object is placed on the marked surface and the scanning is done manually by
moving the scanner around the object at a particular distance. A scanner capable of automatically
scanning objects on a rotating surface was introduced to us. We essentially needed to place
the object on a rotating surface and the scanner/projector will scan the object as it rotates on
the surface.



%------ First Picture -------------
\begin{figure}[h!]
  \center
  \includegraphics[scale=0.31]{./figures/DPC_1.jpg}
\end{figure}

$$
  \frac{X}{f_l} = \frac{B}{Z}
$$


\begin{align*}
  x   & = \text{image coordinate on sensor}       \\
  f_l & = \text{Focal length}                     \\
  B   & = \text{baseline (Distance Laser-sensor)} \\
  Z   & = \text{Object Distance}
\end{align*}


%------ Second Picture -------------
% \begin{figure}[h!]
%   \center
%   \includegraphics[scale=0.5]{./figures/DPC_2.jpg}
% \end{figure}

\end{document}
